\documentclass[11pt,a4paper]{article}
\usepackage[T1]{fontenc}
\usepackage[utf8]{inputenc}
\usepackage{amsmath,amssymb,xcolor}
\usepackage[margin=25mm]{geometry}
\usepackage[colorlinks=true,urlcolor=blue]{hyperref}
\definecolor{nexorablue}{HTML}{2459B5}
\setlength{\parindent}{0pt}
\setlength{\parskip}{7pt}
\setlength{\emergencystretch}{2em}
\renewcommand{\familydefault}{\sfdefault}
\begin{document}

{\small\bfseries\color{nexorablue} NEXORA CPGE}\par
{\small MPSI / NEX-CIR-06 / Énoncé}\par
{\LARGE\bfseries Deux nœuds et un bilan complet}\par
\vspace{0.25em}\hrule\vspace{0.6em}

Une source fixe $V_A=12\,\mathrm V$ par rapport à la masse. Les résistances sont $R_{AB}=2\,\mathrm{k\Omega}$, $R_{BC}=1\,\mathrm{k\Omega}$, $R_{B0}=2\,\mathrm{k\Omega}$ et $R_{C0}=2\,\mathrm{k\Omega}$, où $0$ désigne la masse. Calculer $V_B$, $V_C$ et les quatre courants orientés suivant l’ordre des indices. Comparer la puissance fournie par la source à la somme des puissances dissipées.


\end{document}
